\section{A perturbed cavity reservoir}\label{cavity:section}

We compare systems of sizes $N$ and $N+1$ through a common $N$-spin
system, the \emph{reservoir}. This cavity comparison follows the
Aizenman--Sims--Starr variational scheme~\cite{ASS03} and its
spin-distribution form in Panchenko~\cite[Section~2.2]{Pan13}. We add a sublinear number of bounded perturbations
to obtain identities for its spin trees. The comparison below shows
that these perturbations change the pressure by $o(1)$ and preserve
the site-minus-bond cavity increment.

Until \cref{reduction:section}, assume that $B(\theta)\le B$ and
$|h|\le H$ almost surely and that the law of $\theta$ is invariant under
permutations of its $p$ coordinates. Neither the factorization nor the
positivity assumption is used in the lower bound.

Fix $L\ge2$ and use $r=L-1$ auxiliary levels with
\begin{equation}\label{cavity:grid}
 m_d=d/L\quad(0\le d\le L),\qquad\delta=1/L.
\end{equation}
Let $(c_\nu)_{\nu\ge1}$ be positive with $\sum_\nu c_\nu=1$, and let
$s_N=N^{3/4}$. To the spin weight add independent Poisson families of
single-site factors $\psi_\nu$, with counts of means $c_\nu s_N$ and
independent uniform sites in $[N]$. Each factor uses its own auxiliary
variables, drawn along the hierarchy with product priors, and satisfies
\begin{equation}\label{cavity:factor-bounds}
 1/2\le\psi_\nu\le3/2.
\end{equation}
The total number of factors is $\Pois(s_N)$ and is finite almost surely.
Counts, types, and sites are root variables. The parameters of the factors
are denoted by $\mathbf u$ and are fixed, not quenched random variables,
unless a parameter average is explicitly taken. The factors needed later
are specified in \cref{pert:section}.

Apply the power-mean recursion to the resulting leaf partition function.
Write $P_N^{L,\mathbf u}$ for its expected root logarithm. Comparing
pointwise with the unperturbed spin weight and then applying the recursion
gives
\begin{equation}\label{cavity:negligible}
 |P_N^{L,\mathbf u}-NF_N|\le (\log2)s_N,
\end{equation}
uniformly in the hierarchy, its parameters, and $N$.

\begin{lemma}[Cavity increments]\label[lemma]{cavity:increment}
The increments $\Delta_N^{L,\mathbf u}=P_{N+1}^{L,\mathbf u}-P_N^{L,\mathbf u}$
are uniformly bounded in $N,L,\mathbf u$. There is a reservoir on $N$ spins
with physical interaction count of mean
\begin{equation}\label{cavity:reservoir-mean}
 \rho_N=\alpha(N+1)\left(\frac{N}{N+1}\right)^p,
\end{equation}
the usual fields, and the perturbation counts of means $c_\nu s_N$, such that
\begin{equation}\label{cavity:identity}
 \Delta_N^{L,\mathbf u}
 =\log2+\E\log Y_{0,\rm s}-\E\log Y_{0,\rm b}+o_N(1).
\end{equation}
The error is uniform in $L,\mathbf u$. The insertion recursion uses the
tilted reservoir kernels and independent new root disorder, with terminal
insertions
\begin{align}
 A_{\rm s}(\sigma)&=\frac12\sum_{\eps=\pm1}
 \exp\left(h\eps+\sum_{j=1}^{K_{\rm s}}
 \theta_j(\sigma_{i_{1,j}},\ldots,\sigma_{i_{p-1,j}},\eps)\right),
 \label{cavity:site-insertion}\\
 A_{\rm b}(\sigma)&=\exp\left(\sum_{j=1}^{K_{\rm b}}
 \theta_j(\sigma_{i_{1,j}},\ldots,\sigma_{i_{p,j}})\right),
 \label{cavity:bond-insertion}
\end{align}
where $K_{\rm s}\sim\Pois(\alpha p)$, $K_{\rm b}\sim\Pois(\alpha(p-1))$,
and all displayed indices are independent uniform indices in $[N]$.
\end{lemma}

\begin{proof}
\emph{The physical interactions.}
Split the Poisson interaction process of the $(N+1)$-spin system according
to the number of indices equal to $N+1$. Terms with no new index have mean
$\rho_N$. Terms with exactly one new index have mean
$\alpha p(N/(N+1))^{p-1}=\alpha p+O(N^{-1})$; by permutation invariance we
may place the new input last. Their other indices are independent uniform
indices in $[N]$, and their interaction functions retain the prescribed
law. Terms with at least two new indices have total mean $O(N^{-1})$.
Meanwhile, the $N$-spin system can be obtained from the same reservoir by
adding a Poisson number of old interactions of mean
\[
 \alpha N-\rho_N=\alpha(p-1)+O(N^{-1}).
\]
These quantities are nonnegative for all sufficiently large $N$.

\paragraph{The perturbations.}
The old-site perturbation count in the $(N+1)$-spin system has total mean
$s_{N+1}N/(N+1)$, and the new-site count has mean $s_{N+1}/(N+1)$.
Couple these processes to the reservoir perturbations using common Poisson
points. The expected number of unmatched factors is at most
\[
 |s_{N+1}-s_N|+2s_{N+1}/(N+1)=O(N^{-1/4}).
\]
Deleting or adding a perturbation changes the root logarithm by at most
$\log2$; an interaction changes it by at most $B$. An unused auxiliary
variable can simply be integrated out with its prior, so these comparisons
remain valid when the number of auxiliary variables changes.

\paragraph{The increment.}
We may consequently discard the repeated-new-index interactions, match the
perturbations, and replace the two insertion count means by their limits
at a total expected cost $O(N^{-1/4})$. Summing over the new spin gives
$2A_{\rm s}$, while the old-system correction gives $A_{\rm b}$.
The factor $2$ supplies the term $\log2$ in the increment. The insertion
identity for the power means then proves
\eqref{cavity:identity}. Finally,
\[
 |\log Y_{0,\rm s}|\le H+BK_{\rm s},\qquad
 |\log Y_{0,\rm b}|\le BK_{\rm b}.
\]
Together with the coupling estimate this bounds the increments for large
$N$; increasing the constant covers the remaining sizes. All estimates are
uniform in $L$ and $\mathbf u$.
\end{proof}

In the rest of the lower-bound proof, all spin trees are sampled from this
reservoir using its tilted auxiliary kernels and its terminal Gibbs kernel.
The reservoir retains an independent Poisson interaction process of mean
$O(N)$ and $N$ independent bounded fields. Thus the concentration argument
applies to its slightly modified interaction mean as well.
